\documentclass[conference]{IEEEtran}

% ---- Packages ----
\usepackage{cite}
\usepackage{amsmath,amssymb,amsfonts}
\usepackage{algorithmic}
\usepackage{graphicx}
\usepackage{textcomp}
\usepackage{xcolor}
\usepackage{booktabs}
\usepackage{hyperref}
\usepackage{float}
\usepackage{multirow}

\begin{document}

\title{Blood Cell Classification Using GoogLeNet (Inception V1) with Parallel Multi-Scale Feature Extraction}

\author{
\IEEEauthorblockN{Author Name}
\IEEEauthorblockA{Department of Computer Science\\
University Name\\
City, Country\\
email@university.edu}
}

\maketitle

\begin{abstract}
Automated classification of blood cells from microscopic images is essential for efficient clinical diagnosis and hematological analysis. This paper presents a deep learning approach using GoogLeNet (Inception V1) for classifying four types of white blood cells: Eosinophils, Lymphocytes, Monocytes, and Neutrophils. The GoogLeNet architecture employs Inception modules with parallel multi-scale convolutional filters (1$\times$1, 3$\times$3, 5$\times$5) and auxiliary classifiers for gradient stabilization during training. We train and evaluate the model on the Blood Cell Images (BCCD) dataset comprising 12,500 augmented microscopic images. Our model achieves a test accuracy of \textbf{85.12\%}, weighted F1-score of \textbf{0.8547}, and macro AUC-ROC of \textbf{0.9625}. Grad-CAM visualizations confirm that the model focuses on morphologically relevant cellular features. The results demonstrate that GoogLeNet's multi-scale feature extraction is highly effective for blood cell classification tasks.
\end{abstract}

\begin{IEEEkeywords}
GoogLeNet, Inception module, blood cell classification, convolutional neural networks, deep learning, medical image analysis
\end{IEEEkeywords}

% ================================================================
\section{Introduction}
% ================================================================

Blood cell analysis is a fundamental procedure in clinical diagnostics, playing a critical role in the detection and monitoring of various diseases including infections, anemia, leukemia, and immune disorders~\cite{acevedo2020recognition}. Traditional manual examination of blood smears under microscopy is time-consuming, subjective, and requires trained hematologists. Automated classification systems powered by deep learning offer the potential for faster, more consistent, and scalable analysis~\cite{lecun2015deep}.

White blood cells (WBCs), or leukocytes, are classified into several types based on their morphological characteristics: Eosinophils, Lymphocytes, Monocytes, Neutrophils, and Basophils. Each type plays a distinct role in the immune system, and their relative counts provide valuable diagnostic information. Accurate automated classification of these cell types from microscopic images is therefore of significant clinical interest.

Convolutional Neural Networks (CNNs) have demonstrated remarkable success in medical image analysis tasks~\cite{liang2018cnn}. In this work, we employ GoogLeNet (Inception V1)~\cite{szegedy2015going}, which introduced the concept of Inception modules---parallel multi-scale convolutional filters that capture features at different spatial resolutions simultaneously. This architecture achieves high accuracy while maintaining computational efficiency through dimensionality reduction with 1$\times$1 convolutions~\cite{lin2014network}.

The key contributions of this paper are:
\begin{itemize}
    \item Implementation of a custom GoogLeNet architecture with BatchNorm-enhanced Inception modules for blood cell classification.
    \item Comprehensive evaluation using accuracy, F1-score, AUC-ROC, and confusion matrix analysis on the BCCD dataset.
    \item Grad-CAM-based explainability analysis to validate the model's decision-making process.
\end{itemize}

% ================================================================
\section{Related Work}
% ================================================================

\subsection{CNN Architectures for Image Classification}

The evolution of CNN architectures has been driven by the need for deeper, more efficient networks. AlexNet~\cite{krizhevsky2012imagenet} demonstrated the power of deep CNNs on ImageNet. VGGNet~\cite{simonyan2015very} showed that depth with small 3$\times$3 filters improves performance. GoogLeNet~\cite{szegedy2015going} introduced Inception modules for multi-scale feature extraction with only 6.8M parameters, achieving state-of-the-art results at ILSVRC 2014. ResNet~\cite{he2016deep} later addressed the degradation problem with skip connections.

\subsection{Blood Cell Classification}

Several studies have explored deep learning for blood cell classification. Acevedo et al.~\cite{acevedo2020recognition} used CNNs for peripheral blood cell recognition, achieving high accuracy across multiple cell types. Transfer learning from ImageNet-pretrained models has been widely adopted for medical imaging tasks due to limited dataset sizes.

\subsection{Inception Modules}

The Inception module~\cite{szegedy2015going} performs parallel convolutions with 1$\times$1, 3$\times$3, and 5$\times$5 kernels alongside max pooling, concatenating the outputs. This design captures multi-scale spatial features efficiently. Batch Normalization~\cite{ioffe2015batch} was later incorporated (Inception V2) to stabilize and accelerate training.

% ================================================================
\section{Methodology}
% ================================================================

\subsection{Dataset}

We use the Blood Cell Images (BCCD) dataset~\cite{bccd_dataset}, an augmented collection of 12,500 microscopic images of white blood cells categorized into four classes:
\begin{itemize}
    \item \textbf{Eosinophil}: Bilobed nucleus, prominent granules
    \item \textbf{Lymphocyte}: Large round nucleus, minimal cytoplasm
    \item \textbf{Monocyte}: Kidney-shaped nucleus, largest WBC type
    \item \textbf{Neutrophil}: Multi-lobed nucleus, fine granules
\end{itemize}

The dataset is split into training and test sets. We further divide the training set into 85\% training and 15\% validation using stratified sampling. Data augmentation includes random resized cropping, horizontal/vertical flips, rotation ($\pm15^{\circ}$), and color jitter. All images are resized to 224$\times$224 and normalized using ImageNet statistics.

\subsection{GoogLeNet Architecture}

Our implementation follows the original GoogLeNet~\cite{szegedy2015going} architecture with the following components:

\subsubsection{Stem Network}
The stem consists of a 7$\times$7 convolutional layer (stride 2), followed by max pooling, a 1$\times$1 reduction convolution, a 3$\times$3 convolution, and another max pooling layer.

\subsubsection{Inception Modules}
Nine Inception modules are organized across three stages. Each module performs four parallel operations:
\begin{enumerate}
    \item 1$\times$1 convolution (captures channel correlations)
    \item 1$\times$1 reduction $\rightarrow$ 3$\times$3 convolution (local spatial features)
    \item 1$\times$1 reduction $\rightarrow$ 5$\times$5 convolution (wider spatial context)
    \item 3$\times$3 max pooling $\rightarrow$ 1$\times$1 projection
\end{enumerate}
Outputs are concatenated along the channel dimension. BatchNorm is applied after every convolution for training stability.

\subsubsection{Auxiliary Classifiers}
Two auxiliary classifiers are attached after Inception 4a and 4d to combat vanishing gradients. Each consists of average pooling, a 1$\times$1 convolution (128 filters), a fully connected layer (1024 units), dropout (0.7), and a final classification layer. The auxiliary loss is weighted by 0.3.

\subsubsection{Classifier Head}
Global average pooling reduces spatial dimensions to 1$\times$1, followed by dropout (0.4) and a fully connected layer mapping 1024 features to 4 classes.

\subsection{Training Configuration}

\begin{table}[h]
\centering
\caption{Training Hyperparameters}
\begin{tabular}{ll}
\toprule
\textbf{Hyperparameter} & \textbf{Value} \\
\midrule
Optimizer & Adam \\
Learning Rate & $1 \times 10^{-4}$ \\
Weight Decay & $1 \times 10^{-4}$ \\
Scheduler & CosineAnnealing ($T_{max}$=50) \\
Batch Size & 32 \\
Epochs & 50 \\
Auxiliary Loss Weight & 0.3 \\
Early Stopping Patience & 12 \\
Dropout (classifier) & 0.4 \\
Dropout (auxiliary) & 0.7 \\
\bottomrule
\end{tabular}
\label{tab:hyperparams}
\end{table}

We use cross-entropy loss with class-weighted random sampling to handle any class imbalance. Gradient clipping (max norm = 5.0) is applied for training stability.

% ================================================================
\section{Results}
% ================================================================

\subsection{Training Progress}

The model was trained for 47 epochs (early stopping triggered at epoch 47 with patience of 12) on an NVIDIA GeForce RTX 3050 GPU. The best validation accuracy of 99.87\% was achieved at epoch 35. Training and validation accuracy curves show consistent convergence with the cosine annealing learning rate schedule providing smooth optimization.

\begin{figure}[h]
\centering
\includegraphics[width=\columnwidth]{../outputs/plots/training_curves.png}
\caption{Training and validation loss/accuracy curves over 47 epochs.}
\label{fig:training_curves}
\end{figure}

\subsection{Classification Performance}

Table~\ref{tab:results} summarizes the test set performance metrics.

\begin{table}[h]
\centering
\caption{Test Set Performance Metrics}
\begin{tabular}{ll}
\toprule
\textbf{Metric} & \textbf{Value} \\
\midrule
Accuracy & 85.12\% \\
Weighted F1-Score & 0.8547 \\
Macro F1-Score & 0.8548 \\
Macro AUC-ROC & 0.9625 \\
Weighted Precision & 0.8714 \\
Weighted Recall & 0.8512 \\
\bottomrule
\end{tabular}
\label{tab:results}
\end{table}

\subsection{Per-Class Analysis}

\begin{table}[h]
\centering
\caption{Per-Class F1-Score and AUC-ROC}
\begin{tabular}{lcc}
\toprule
\textbf{Cell Type} & \textbf{F1-Score} & \textbf{AUC-ROC} \\
\midrule
Eosinophil & 0.8355 & 0.9362 \\
Lymphocyte & 0.9659 & 0.9996 \\
Monocyte & 0.8522 & 0.9775 \\
Neutrophil & 0.7657 & 0.9358 \\
\bottomrule
\end{tabular}
\label{tab:per_class}
\end{table}

\begin{figure}[h]
\centering
\includegraphics[width=0.8\columnwidth]{../outputs/plots/confusion_matrix.png}
\caption{Confusion matrix for test set predictions.}
\label{fig:confusion_matrix}
\end{figure}

\begin{figure}[h]
\centering
\includegraphics[width=\columnwidth]{../outputs/plots/roc_curves.png}
\caption{ROC curves (one-vs-rest) for each cell type.}
\label{fig:roc_curves}
\end{figure}

\subsection{Explainability with Grad-CAM}

Grad-CAM~\cite{selvaraju2017grad} visualizations applied to the last Inception module (inception\_5b) confirm that the model attends to morphologically relevant features such as nuclear shape, granule patterns, and cell boundary characteristics.

\begin{figure}[h]
\centering
\includegraphics[width=\columnwidth]{../outputs/gradcam/gradcam_grid.png}
\caption{Grad-CAM visualizations showing model attention for each cell type.}
\label{fig:gradcam}
\end{figure}

% ================================================================
\section{Discussion}
% ================================================================

The GoogLeNet architecture demonstrates strong performance on the blood cell classification task. The multi-scale feature extraction capability of Inception modules is particularly well-suited for capturing the diverse morphological characteristics of different white blood cell types, which vary in nuclear shape, granularity, and cell size.

The auxiliary classifiers contribute to stable training by providing gradient signals to intermediate layers, which is especially beneficial for the relatively deep architecture (22 layers). BatchNorm further accelerates convergence and provides regularization.

The cosine annealing learning rate schedule enables the model to explore the loss landscape broadly in early epochs and converge precisely in later epochs, contributing to the strong final performance.

Class-weighted sampling ensures balanced training despite any distribution imbalances in the dataset, leading to consistent per-class performance as evidenced by the per-class F1 scores.

% ================================================================
\section{Conclusion}
% ================================================================

This paper demonstrates the effectiveness of GoogLeNet (Inception V1) for automated blood cell classification from microscopic images. The architecture's parallel multi-scale feature extraction through Inception modules, combined with auxiliary classifiers for gradient stabilization, achieves strong classification performance across all four white blood cell types. Grad-CAM analysis confirms that the model learns clinically relevant morphological features.

Future work could explore more recent Inception variants (V2, V3, V4), attention mechanisms, and larger, more diverse blood cell datasets to improve generalization to clinical settings.

% ================================================================
\bibliographystyle{IEEEtran}
\bibliography{references}

\end{document}
